The quantification of uncertainty has emerged as a significant area of research across various machine learning domains, including natural language processing (NLP). 
However, previous NLP studies have predominantly addressed the associated UQ challenges similarly to classification or regression methodologies~\citep{desai-durrett-2020-calibration,jiang-etal-2021-know,kamath-etal-2020-selective,wang-etal-2022-uncertainty,xiong2023llms}. 
For instance, \citet{kamath-etal-2020-selective} examines the selective question-answering task as a multiple-choice problem, reducing it to a de facto classification task rather than directly engaging with free-form generation. 
As recently argued in \cite{kuhn2023semantic}, such approaches enable the application of UQ measures akin to those employed in more extensively researched classification or regression contexts, but overlook the generative aspects and distinct challenges of NLG.

Recently, some research started to study uncertainty quantification for NLG.
One line of research involves asking the LLM itself for its confidence, with or without additional fine-tuning~\citep{kadavath2022language,Lin2022TeachingMT, DBLP:journals/corr/abs-2012-14983,chen2023quantifying}.
Apart from being expensive, such approaches can be hard to generalize due to opaque training details or differences between LLMs~\citep{kuhn2023semantic}. 
The work most relevant to ours is~\cite{kuhn2023semantic}, which proposes to compute the ``semantic entropy'' by considering the equivalence relationships amongst generated answers, and requires no training.
Nonetheless, it still requires access to the token-level numerical output of the LLM, which is not always available.

As discussed, one of the most pertinent applications of uncertainty quantification in NLG involves the development of methods for selective NLG (or, NLG with rejection). 
This emerging field has limited research to date, but shares close ties with \textit{classification with rejection}.
Both tasks can be viewed as determining when to trust a model, whether it is a classifier or an LLM. 
Numerous classification with rejection methods emphasize the identification of a reliable confidence score (some of which are jointly trained with the classifier)~\citep{Corbiere2019AddressingConfidence,Fumera2000RejectThresholds,Geifman2017SelectiveNetworks,Jiang2018ToClassifier}, which is often not only dependent on the input but also on the prediction. 
As existing uncertainty quantification research for NLG primarily focuses on input uncertainty~\citep{kuhn2023semantic,malinin2021uncertainty}, it overlooks the crucial aspect of confidence, which is essential in deciding when to trust an LLM's response (see \cref{sec:background:uvsc} for more discussion). 
Recent works have explored selective classification in NLP tasks~\citep{varshney-etal-2022-investigating,varshney-etal-2022-towards}.
However, the distinct generative nature of NLG precludes the direct adaptation of confidence measures from the classification with rejection literature. 
This paper serves as a step to bridge this gap and enhance the effectiveness of uncertainty quantification in NLG. 


It is worth noting that the issue of LLMs being overconfident \citep{mielke-etal-2022-reducing,si-etal-2022-examining,xiong2023llms} is orthogonal to our work, as we evaluate measures basing how they \textit{rank} different samples - such measures may then be calibrated by distribution-free uncertainty quantification methods like in \citet{schuster2022confident}\footnote{
See Appendix for results on calibrated confidence measures.
}.
\citet{giulianelli2023comes} also provides an interesting exploration of the inherent uncertainty in human responses for many NLG tasks, in a black-box manner.
Finally, carefully designed prompts have been proposed to improve the quality of the generated responses in general~\citep{Zhou2023NavigatingTG,si2023prompting,wei2022chain}.
Orthogonal to UQ but related to selective NLG, \cite{varshney-baral-2023-post} focuses on reattempting rejected samples, with the help of an auxiliary model trained on an additional dataset that predicts the correctness of the generation.
This paper focuses on providing quantitative uncertainty/confidence measures, which can be used to identify high-quality generations. 
